% ============================================================
% B-checks.tex  -- GENERATED by gen_checks_B.py; do not edit by hand.
% The four check circuits of the two-ancilla layout in program order, from
% LRB-experiment-data-slurm/comparisons/two-ancilla-checks/hw-layouts/encoders/layoutB_round_ops.txt
% (wire 5 = N, wire 6 = E). Layout: rows are the two halves of the round
% (G2, G1 on top; G4, G3 below), columns are the ancillas (N left, E right).
% Checked gate by gate, and by an exact statevector run of each block, by
% verify_checks_B.py.
% ============================================================
\documentclass[border=1pt]{standalone}
\usepackage{amsmath,amssymb}
\usepackage{xcolor}
\usepackage{quantikz}
\definecolor{Zcol}{HTML}{FF1744}
\definecolor{Xcol}{HTML}{00B7FF}
\colorlet{XText}{Xcol!80!black}
\colorlet{ZText}{Zcol!85!black}
\newcommand{\rowsep}{0.62cm}   % wire spacing before scaling
\newcommand{\colsep}{0.20cm}   % gate column spacing before scaling
\newcommand{\titlegap}{2pt}
\newlength{\figw}\setlength{\figw}{16.0cm}   % final width
\newlength{\colgap}\setlength{\colgap}{0.55cm}
\newlength{\rowgap}\setlength{\rowgap}{0.30cm}
\newlength{\leftw}
\newsavebox{\boxone}\newsavebox{\boxtwo}\newsavebox{\boxthree}\newsavebox{\boxfour}
\newcommand{\Panel}[2]{{\small\bfseries(#1)}\hspace{0.45em}{\footnotesize #2}}
\begin{document}
\sbox{\boxone}{%
% G2: [('CNOT', 2, 5), ('CNOT_INV', 1, 4), ('CNOT_INV', 4, 5), ('CNOT', 1, 4)]
\begin{quantikz}[row sep={\rowsep,between origins}, column sep=\colsep,
  execute at end picture={\path ([yshift=0.42cm]\tikzcdmatrixname-1-1.center);}]
      \lstick{$q_1$} & \qw & \ctrl{2} & \qw & \ctrl{2} & \qw \\
      \lstick{$q_2$} & \ctrl{2} & \qw & \qw & \qw & \qw \\
      \lstick{$q_4$} & \qw & \gate{X^{-1}} & \ctrl{1} & \gate{X} & \qw \\
      \lstick{$\mathrm{N}: \ket{0}$} & \gate{X} & \qw & \gate{X^{-1}} & \qw & \meter{}
\end{quantikz}%
}%
\sbox{\boxtwo}{%
% G1: [('CNOT', 2, 6), ('CNOT_INV', 0, 1), ('CNOT', 1, 2), ('CNOT', 3, 2), ('CNOT_INV', 2, 6), ('CNOT_INV', 3, 2), ('CNOT_INV', 1, 2), ('CNOT', 0, 1)]
\begin{quantikz}[row sep={\rowsep,between origins}, column sep=\colsep,
  execute at end picture={\path ([yshift=0.42cm]\tikzcdmatrixname-1-1.center);}]
      \lstick{$q_0$} & \qw & \ctrl{1} & \qw & \qw & \qw & \qw & \qw & \ctrl{1} & \qw \\
      \lstick{$q_1$} & \qw & \gate{X^{-1}} & \ctrl{1} & \qw & \qw & \qw & \ctrl{1} & \gate{X} & \qw \\
      \lstick{$q_2$} & \ctrl{2} & \qw & \gate{X} & \gate{X} & \ctrl{2} & \gate{X^{-1}} & \gate{X^{-1}} & \qw & \qw \\
      \lstick{$q_3$} & \qw & \qw & \qw & \ctrl{-1} & \qw & \ctrl{-1} & \qw & \qw & \qw \\
      \lstick{$\mathrm{E}: \ket{0}$} & \gate{X} & \qw & \qw & \qw & \gate{X^{-1}} & \qw & \qw & \qw & \meter{}
\end{quantikz}%
}%
\sbox{\boxthree}{%
% G4: [('H', 5), ('CNOT', 5, 2), ('CNOT_INV', 4, 1), ('CNOT_INV', 5, 4), ('CNOT', 2, 3), ('CNOT_INV', 5, 2), ('CNOT_INV', 2, 3), ('CNOT', 4, 1), ('H_INV', 5)]
\begin{quantikz}[row sep={\rowsep,between origins}, column sep=\colsep,
  execute at end picture={\path ([yshift=0.42cm]\tikzcdmatrixname-1-1.center);}]
      \lstick{$q_1$} & \qw & \qw & \gate{X^{-1}} & \qw & \qw & \qw & \qw & \gate{X} & \qw & \qw \\
      \lstick{$q_2$} & \qw & \gate{X} & \qw & \qw & \ctrl{1} & \gate{X^{-1}} & \ctrl{1} & \qw & \qw & \qw \\
      \lstick{$q_3$} & \qw & \qw & \qw & \qw & \gate{X} & \qw & \gate{X^{-1}} & \qw & \qw & \qw \\
      \lstick{$q_4$} & \qw & \qw & \ctrl{-3} & \gate{X^{-1}} & \qw & \qw & \qw & \ctrl{-3} & \qw & \qw \\
      \lstick{$\mathrm{N}$} & \gate{H} & \ctrl{-3} & \qw & \ctrl{-1} & \qw & \ctrl{-3} & \qw & \qw & \gate{H^{-1}} & \meter{}
\end{quantikz}%
}%
\sbox{\boxfour}{%
% G3: [('H', 6), ('CNOT_INV', 1, 0), ('CNOT', 2, 1), ('CNOT_INV', 6, 2), ('CNOT_INV', 2, 1), ('CNOT', 1, 0), ('H_INV', 6)]
\begin{quantikz}[row sep={\rowsep,between origins}, column sep=\colsep,
  execute at end picture={\path ([yshift=0.42cm]\tikzcdmatrixname-1-1.center);}]
      \lstick{$q_0$} & \qw & \gate{X^{-1}} & \qw & \qw & \qw & \gate{X} & \qw & \qw \\
      \lstick{$q_1$} & \qw & \ctrl{-1} & \gate{X} & \qw & \gate{X^{-1}} & \ctrl{-1} & \qw & \qw \\
      \lstick{$q_2$} & \qw & \qw & \ctrl{-1} & \gate{X^{-1}} & \ctrl{-1} & \qw & \qw & \qw \\
      \lstick{$\mathrm{E}$} & \gate{H} & \qw & \qw & \ctrl{-1} & \qw & \qw & \gate{H^{-1}} & \meter{}
\end{quantikz}%
}%
% one scale for all four blocks: the wider block of each column sets the
% column width, and the two columns fill \figw minus \colgap
\pgfmathsetmacro{\wleft}{max(\wd\boxone,\wd\boxthree)}%
\pgfmathsetmacro{\wright}{max(\wd\boxtwo,\wd\boxfour)}%
\pgfmathsetmacro{\circscale}{(\figw-\colgap)/(\wleft+\wright)}%
\setlength{\leftw}{\circscale\dimexpr\wleft pt\relax}%
\typeout{CIRCSCALE \circscale\space L=\wleft\space R=\wright\space leftw=\the\leftw}%
\begin{minipage}{\figw}
\begin{minipage}[t]{\leftw}
\Panel{a}{\textcolor{ZText}{$G_2 = Z_1 Z_2 Z_4^{-1}$}, read by N}\\[\titlegap]
\scalebox{\circscale}{\usebox{\boxone}}
\end{minipage}\hspace{\colgap}%
\begin{minipage}[t]{\dimexpr\figw-\leftw-\colgap\relax}
\Panel{b}{\textcolor{ZText}{$G_1 = Z_0 Z_1^{-1} Z_3^{-1}$}, read by E}\\[\titlegap]
\scalebox{\circscale}{\usebox{\boxtwo}}
\end{minipage}

\vspace{\rowgap}

\begin{minipage}[t]{\leftw}
\Panel{c}{\textcolor{XText}{$G_4 = X_1^{-1} X_3 X_4^{-1}$}, read by N}\\[\titlegap]
\scalebox{\circscale}{\usebox{\boxthree}}
\end{minipage}\hspace{\colgap}%
\begin{minipage}[t]{\dimexpr\figw-\leftw-\colgap\relax}
\Panel{d}{\textcolor{XText}{$G_3 = X_0 X_1 X_2^{-1}$}, read by E}\\[\titlegap]
\scalebox{\circscale}{\usebox{\boxfour}}
\end{minipage}
\end{minipage}
\end{document}
